\vspace{-2pt}
\section{Related work}
\label{sec:related}
\vspace{-4pt}

% \js{Move this sec. to the end? It could help introduce the trials and dataset early in the paper}

Energy-aware management of cellular access networks has been extensively studied. A seminal work by Ajmone et al. \cite{Ajmone2009} characterized energy savings by reducing active cells during low-traffic periods. However, this and subsequent studies often assume uniform traffic across cells and complete offloading to a predefined set of fully overlapping neighboring cells. Further research has focused on optimal sleeping policies derived specifically for non-bursty traffic. These studies reveal a two-threshold policy with a wait-and-see feature, where the server waits to see if there are additional arrivals before switching modes \cite{two_thresholds_traffic}. Despite the lack of theoretical results for multiple servers, these findings align with current carrier shutdown solutions that implement distinct conditions for shutdown and reactivation \cite{data_driven_energy23}.

While numerous studies have explored enhancing energy efficiency in telecommunication networks through \textit{energy consumption measurement} and \textit{policies for optimizing energy consumption}\cite{survey_5g_energy_efficiency, survey2011}, a gap remains in understanding the real-world impact of these policies. To sensibly assess the performance of different policies, it is necessary to conduct regional trials. This approach would provide a real evaluation of the energy gains or savings achieved by each policy.

\textbf{Measuring energy consumption.} The energy consumption of mobile networks is a major concern for the telecom industry. To evaluate the energy-related performance of the RAN, the industry has defined various measurement methods and metrics across different network levels—network \cite{ETSI2020a}, site \cite{ITU2016}, base station (BS) \cite{ETSI2020b}, and user equipment \cite{3GPP2019}—as well as for different scenarios such as dense urban, urban, and rural coverage \cite{ETSI2020a}, and services including enhanced mobile broadband, ultra-reliable low-latency communications, and massive machine-type communications \cite{3GPP2021}. Piovesan et al. \cite{Piovesan2022} propose a model for characterizing the power consumption of 5G multi-carrier BSs, built on extensive data collection. Additionally, Lopes et al. \cite{survey_5g_energy_efficiency} conducted a survey outlining the primary energy efficiency technologies provided by 3GPP NR, detailing their benefits and challenges. The community recognizes the importance of sustainability, with efforts underway to quantify the carbon footprint of communication networks \cite{IEA2024} and optimize network routing \cite{zilberman23}, alongside general efforts for computer systems \cite{CarbonScaler}. In this study, we present utilization and performance data collected at various times throughout the day. This information is instrumental in assessing whether base station utilization can benefit from temporal flexibility, a crucial aspect in achieving greater carbon efficiency.

\textbf{Optimizing energy saving policies.} Cell sleep policies involve the dynamic shutdown of cells during low-traffic periods and their subsequent reactivation as traffic increases. Accurately predicting traffic loads in the near future is crucial for determining the appropriate times to switch base stations on and off. Marsan et al. \cite{Marsan2009}\cite{Marsan2012} evaluated various fixed switching-off schemes in cellular networks, finding that optimal configurations vary with traffic profiles. They later showed that cooperative strategies achieve significant savings among operators for user roaming and energy savings \cite{Marsan2010}. Donevski et al \cite{Donevski2019} propose using dense Neural Network and Recurrent Neural Network paradigms, although their approach lacks empirical validation through real-world trials. Moreover, Domenico et al. \cite{user_transfer_5g} focus on the carrier shutdown approach, enabling base stations to autonomously power down during low-traffic periods by transferring their load to neighboring active base stations. Other researchers have employed reinforcing learning methods\cite{rl2020}\cite{rl2021}, but such approaches have been tested only in simulated environments. It is crucial to acknowledge that such optimization involves the sudden removal and addition of carriers to the network. Therefore, there is a need to develop automated and precise methods for generating configuration parameters for newly added carriers in cellular networks (\ie planning), as suggested by \cite{auric}. These types of solutions are orthogonal to previously mentioned works and to the energy-saving strategies analyzed in this paper.